% preamble.tex — Shared LaTeX preamble for Instructor's Manual conversion
% Applied PDE with Fourier Series and BVP, 4th Ed. — Richard Haberman

\documentclass[10pt]{book}

% TOC: show only chapters (matching original manual's clean 14-entry TOC)
\setcounter{tocdepth}{0}

%% ── Page Geometry ──────────────────────────────────────────────
% US Letter with generous margins matching scanned source
\usepackage[
  letterpaper,
  top=1in,
  bottom=1.2in,
  left=1.1in,
  right=1.1in,
  headheight=14pt,
  headsep=0.3in,
  footskip=0.4in
]{geometry}

%% ── Encoding & Fonts ──────────────────────────────────────────
\usepackage[T1]{fontenc}
\usepackage[utf8]{inputenc}
\usepackage{lmodern}
\usepackage{microtype}

%% ── Math ───────────────────────────────────────────────────────
\usepackage{amsmath}
\usepackage{amssymb}
\usepackage{amsthm}
\usepackage{mathtools}
\usepackage{bm}

%% ── Solution Environment ────────────────────────────────────────
% Solutions are labeled by original exercise number: "1.4.1 (a)"
\usepackage{enumitem}

% Each solution block: \begin{solution}{1.4.1} ... \end{solution}
% Content uses (a), (b), etc. as sub-items when needed
\newenvironment{solution}[1]{%
  \par\medskip\noindent\textbf{#1}\quad\ignorespaces
}{\par\medskip}

% For multi-part solutions: \part{a}, \part{b}, etc.
\newcommand{\solpart}[1]{\par\noindent(#1)\quad}

%% ── Section Headings ────────────────────────────────────────────
% The manual uses "Section X.Y" as sub-headings within each chapter
\usepackage{titlesec}
\titleformat{\section}[block]{\bfseries\large}{}{0pt}{Section }
\titlespacing{\section}{0pt}{1.5em}{0.8em}

%% ── Figures (minimal — solutions manual has very few) ──────────
\usepackage{graphicx}
\usepackage{float}

%% ── Headers & Footers ──────────────────────────────────────────
\usepackage{fancyhdr}
\pagestyle{fancy}
\fancyhf{}
\fancyhead[LE]{\thepage}
\fancyhead[RE]{\leftmark}
\fancyhead[RO]{\thepage}
\fancyhead[LO]{\rightmark}
\renewcommand{\headrulewidth}{0pt}

\renewcommand{\chaptermark}[1]{\markboth{Chapter \thechapter.\ #1}{}}
\renewcommand{\sectionmark}[1]{\markright{\thesection\ #1}}

\fancypagestyle{plain}{%
  \fancyhf{}%
  \fancyfoot[C]{\thepage}%
  \renewcommand{\headrulewidth}{0pt}%
}

%% ── Cross-references ──────────────────────────────────────────
\usepackage{hyperref}
\hypersetup{
  colorlinks=false,
  pdfborder={0 0 0},
  pdftitle={Instructor's Manual: Applied PDE with Fourier Series and BVP},
  pdfauthor={Richard Haberman}
}

%% ── Custom Commands ────────────────────────────────────────────

% Partial derivatives
\newcommand{\pd}[2]{\frac{\partial #1}{\partial #2}}
\newcommand{\pdd}[2]{\frac{\partial^2 #1}{\partial #2^2}}
\newcommand{\pddm}[3]{\frac{\partial^2 #1}{\partial #2 \partial #3}}

% Ordinary derivatives
\newcommand{\od}[2]{\frac{d #1}{d #2}}
\newcommand{\odd}[2]{\frac{d^2 #1}{d #2^2}}

% Evaluation bar
\newcommand{\evalat}[2]{\left. #1 \right|_{#2}}

% Common operators
\DeclareMathOperator{\erfc}{erfc}
\DeclareMathOperator{\sinc}{sinc}
\DeclareMathOperator{\sgn}{sgn}
\DeclareMathOperator{\Real}{Re}
\DeclareMathOperator{\Imag}{Im}

% Bold vectors
\renewcommand{\vec}[1]{\mathbf{#1}}

% Laplacian
\newcommand{\laplacian}{\nabla^2}

% Differentials
\newcommand{\dx}{\,dx}
\newcommand{\dy}{\,dy}
\newcommand{\dt}{\,dt}
\newcommand{\ds}{\,ds}
\newcommand{\dA}{\,dA}
\newcommand{\dV}{\,dV}
\newcommand{\dS}{\,dS}
\newcommand{\dr}{\,dr}
\newcommand{\dtheta}{\,d\theta}
